% !TeX root = ../main.tex

\chapter{总结与展望}

本文围绕单、双变量含参函数的“任意性”与“存在性”问题，对其求解方法和教学策略进行了系统研究。

在理论研究方面，我们从一个具体的函数例子出发，通过改变函数个数、变量个数以及变换量词组合等方法，自然地引出了各类相关题型，建立了完整的题目类型体系。

在教学方法上，提出了“案例引导—动态演示—分层探究”的教学模式。以工厂营收问题为现实情境，借助GeoGebra软件实现参数变化的可视化呈现，帮助学生建立直观理解。采用教师示范（单函数问题）与学生探究（双函数问题）相结合的方式，既能保证教学效率，又能促进学生深度参与。

本研究的主要特色在于：一是通过典型例题的变式设计，将各类题型有机联系起来；二是采用“问题阶梯”教学法，设计一个实际背景下的问题链，串联起所有题型的教学，体现了不同题型之间的差异和联系；三是借助动态几何软件，使抽象的数学概念变得形象直观；四是采用讲练结合的方式，兼顾了教学效率和学习效果；五是采用了教师引导教学和学生自主探究结合的方式，培养学生的自主学习能力。当然，这项研究还存在一些不足，比如实际案例可以更丰富，学生的个体差异需要更多关注等。希望未来能有更多学者继续完善教学方法，开发更多实用的教学资源，让更多学生能够更好地掌握这类重要的数学问题。相信通过不断改进，我们能够找到更有效的教学途径，帮助学生提升数学思维能力和问题解决能力。